\section{Run temperature export}

The command \command{bdx_run_temperature_export} retrieves the configured Raspberry temperature PVs from EPICS Archiver Appliance over the interval defined by a DAQ run information file. It can produce a JSON summary and, when the optional export dependencies are installed, a ROOT file for downstream DAQ analysis.

ROOT support is intentionally optional so IOC-only hosts and the Raspberry Pi do not acquire NumPy or uproot. Install it on the DAQ or analysis host with:

\begin{lstlisting}[style=shell]
python -m pip install -e '.[export]'
\end{lstlisting}

The default configuration is \file{config/run_temperature_export.json}. It defines the DAQ run root, run-directory and information-file patterns, the Europe/Rome run time zone, the Archiver retrieval endpoint and timeout, the output path, histogram bin width, missing-alarm sentinel, and the temperature PV list.

Create the final ROOT file with:

\begin{lstlisting}[style=shell]
bdx_run_temperature_export 260714 --write-root
\end{lstlisting}

The output is written atomically to \file{<run_directory>/OTHER/SlowControl_<run_id>.root}. The \file{OTHER} directory must already exist. An existing output is replaced only when \code{--overwrite} is supplied.

The ROOT output contains:

\begin{itemize}
  \item a \code{temperature_samples} TTree with normalized timestamps, values, status, and severity;
  \item mean and sample-count TH1D histograms for each configured sensor;
  \item a \code{metadata} TObjString containing UTF-8 JSON describing the run interval, configuration, warnings, and output provenance.
\end{itemize}

Missing Archiver status or severity fields use the configured \code{missing_alarm_value} sentinel, currently \code{-1}; the sentinel is also recorded in the metadata. Retrieval is timestamp-based, therefore the DAQ host, IOC hosts, and Archiver host must use synchronized clocks and the run information file must be interpreted in the configured time zone.

When modifying this workflow, preserve atomic output semantics, explicit overwrite protection, strict Archiver-response validation, and tests covering run-file parsing, time-zone conversion, empty intervals, malformed samples, ROOT structure, and metadata consistency.
